% PAGES: 278-278
% Closes the \begin{answerx} opened at the end of sec09_c1_11.tex (PDF page
% 277): the printed closing square appears after "...respectively." below.
%
% This is chunk c1's LAST file. Cross-checked against chunk c2's
% sec09_c2_1.tex (delivered before this file, since c2 started before any c1
% file existed and rendered PDF page 278 itself to judge the boundary): c2's
% note states that page 278 ends with the $w_{min}$ derivation closing at
% "$= (1-w_{min,cash})w_T,$" tagged (9.39), "a numbered multi-line display
% that is fully closed on that page, not carried over," and that PDF page 279
% (c2's own first page) opens a fresh sentence, "since, from (9.34), it
% follows that", with its own self-contained unnumbered display. My own
% independent reading of page 278 (rendered directly via CropBox, 600 dpi)
% agrees exactly: nothing else is printed below eq. (9.39) on this page. The
% \begin{align} below is therefore closed with \end{align} in this file, NOT
% left open -- the trailing comma after (9.39) belongs to a sentence that
% continues in prose (not a further display-math line) into c2's
% sec09_c2_1.tex, which already picks it up correctly. No environment is left
% open at the end of chunk c1.

and the weight of the cash position is
\[
w_{min,cash} \;=\; 1-\bfone^tw_{min} \;=\; -0.3356.
\]

Thus, the \$100 million minimum variance portfolio with expected return equal to
2.25\% is obtained by borrowing \$33.56 million and investing \$21.20 million in asset
1, \$48.88 million in asset 2, \$35.40 million in asset 3, and \$28.08 million in asset
4.

From (9.35), we find that the standard deviation of the return of the minimum
variance portfolio is $\sigma_{min} = 0.1949$, i.e., 19.49\%, which is smaller than the
standard deviation of the returns of each of the four assets, which are $\sqrt{0.09} =
0.3 = 30\%$, $\sqrt{0.0625} = 0.25 = 25\%$, $\sqrt{0.1225} = 0.35 = 35\%$,
$\sqrt{0.0576} = 0.24 = 24\%$, respectively. \hfill$\square$
\end{answerx}

\subsection{Minimum variance portfolios and the tangency portfolio}

It is important to note the connection between the minimum variance portfolio and a
special portfolio called the tangency portfolio: in a nutshell, any minimum variance
portfolio consists of a cash position that is determined to match the required
expected return of the portfolio, with the rest of the portfolio invested in the
tangency portfolio.

To see this, and implicitly introduce the tangency portfolio, note that the asset
weights vector $w_{min}$ given by (9.33) is a scalar multiple of the vector
$\Sigma_R^{-1}\bar\mu$, i.e., $w_{min} = c\,\Sigma_R^{-1}\bar\mu$, where $c =
\frac{\mu_P-r_f}{\bar\mu^t\Sigma_R^{-1}\bar\mu}$. Denote by
\begin{equation}
w_T \;=\; \frac{1}{\bfone^t\Sigma_R^{-1}\bar\mu}\Sigma_R^{-1}\bar\mu
\end{equation}
the asset weights vector in a portfolio with no cash, i.e., with $\bfone^tw_T = 1$,
and with the $n$ assets allocated in the same proportions as the $\Sigma_R^{-1}\bar\mu$
vector.

This portfolio is called the \textit{tangency portfolio}. Note that another way of
identifying the tangency portfolio is as the portfolio with the highest Sharpe ratio of
all portfolios fully invested in $n$ assets. The tangency portfolio also plays an
important role in the Capital Asset Pricing Model (CAPM), where it is called the market
portfolio.

From (9.36), we find that
\[
\bar\mu^tw_T \;=\; \frac{\bar\mu^t\Sigma_R^{-1}\bar\mu}{\bfone^t\Sigma_R^{-1}\bar\mu},
\]
and therefore
\begin{equation}
\frac{\bfone^t\Sigma_R^{-1}\bar\mu}{\bar\mu^t\Sigma_R^{-1}\bar\mu} \;=\;
\frac{1}{\bar\mu^tw_T}.
\end{equation}

From (9.34) and (9.37), we obtain that the weight $w_{min,cash}$ of the cash position
in the minimum variance portfolio can be written in terms of $w_T$ as follows:
\begin{equation}
w_{min,cash} \;=\; 1 - \frac{\mu_P-r_f}{\bar\mu^tw_T}.
\end{equation}

Moreover, from (9.33) and (9.36), we find that
\begin{align}
w_{min} &= \frac{\mu_P-r_f}{\bar\mu^t\Sigma_R^{-1}\bar\mu} \cdot \Sigma_R^{-1}\bar\mu
\notag\\
&= (\mu_P-r_f)\frac{\bfone^t\Sigma_R^{-1}\bar\mu}{\bar\mu^t\Sigma_R^{-1}\bar\mu}
\cdot \frac{1}{\bfone^t\Sigma_R^{-1}\bar\mu}\Sigma_R^{-1}\bar\mu \notag\\
&= (1-w_{min,cash})\,w_T,
\end{align}
% Chunk c1 ends here (PDF page 278 / folio 262 is fully transcribed). The
% \end{align} above closes cleanly on this page -- confirmed against chunk
% c2's sec09_c2_1.tex, whose opening line "since, from (9.34), it follows
% that" is the direct prose continuation of the sentence ending in the
% trailing comma after (9.39), not a further aligned display line. c2's file
% already exists and needs no changes.
